\documentclass[10pt,a4paper]{report}
\usepackage[utf8]{inputenc}
\usepackage{amsmath}
\usepackage{amsfonts}
\usepackage{amssymb}
\usepackage{amsthm}
\usepackage{hyperref}

\usepackage{multicol}
\usepackage{fancyhdr}
\usepackage[inline]{enumitem}
\usepackage{tikz}
\usepackage{tikz-cd}
\usetikzlibrary{calc}
\usetikzlibrary{shapes.geometric}
\usepackage[margin=0.5in]{geometry}
\usepackage{xcolor}
\usepackage{ esint }

\hypersetup{
    colorlinks=true,
    linkcolor=blue,
    filecolor=magenta,      
    urlcolor=cyan,
    pdftitle={Tensors},
    pdfpagemode=FullScreen,
    }

%\urlstyle{same}

\newcommand{\CLASSNAME}{Math 5230 -- Partial Differential Equations}
\newcommand{\STUDENTNAME}{Paul Carmody}
\newcommand{\ASSIGNMENT}{Homework \#3 }
\newcommand{\DUEDATE}{October 15, 2025}
\newcommand{\SEMESTER}{Fall 2025}
\newcommand{\SCHEDULE}{MW 2:00 -- 3:15}
\newcommand{\ROOM}{Crawford 332}

\newcommand{\MMN}{M_{m\times n}}
\newcommand{\FF}{\mathcal{F}}

\pagestyle{fancy}
\fancyhf{}
\chead{ \fancyplain{}{\CLASSNAME} }
%\chead{ \fancyplain{}{\STUDENTNAME} }
\rhead{\thepage}
\input{../MyMacros}

\newcommand{\RED}[1]{\textcolor{red}{#1}}
\newcommand{\BLUE}[1]{\textcolor{blue}{#1}}

\begin{document}

\begin{center}
	\Large{\CLASSNAME -- \SEMESTER} \\
	\large{ w/Professor Snelson}
\end{center}
\begin{center}
	\STUDENTNAME \\
	\ASSIGNMENT -- \DUEDATE\\
\end{center} 

Please show your work and justify all answers.

\begin{enumerate}
\item Show by direct calcuation that the fundamental solution 
\begin{align*}
	\Phi(x,t) = \frac{1}{(4\pi t)^{n/2}}e^{-|x|^2/4t}
\end{align*}solves the heat equation $\Phi_t-=\Delta \Phi = 0$ for any $ x \in \R^n$ and $t>0$.

\item Assume $u$ is a solution fo the heat equation on a bounded domain:
\begin{align*}
	u_t-\Delta u=0,\, t>0, x\in U.
\end{align*}Show that if $u$ satisfies the Dirichlet boundary condiotn $u=0$ on $\partial U$ OR the Neumann boundary condition $\PART{u}{\nu}=0$ on $\partial U$, then the quantity $\int_U u^2(x,t)\,dx$ is decreasing s a function of time.

\item Consider the heat equation on a bounded domain with Newmann boundary conditions on $\partial U$:
\begin{align*}
	\BINDEF{u_t-\Delta u=0, & x \in U, t>0}{\PART{u}{\nu}=0,& x \in \partial U}
\end{align*}Show that the quantity $\in_U\BARS{u(x,t)}^2\,dx$ is decreasing as  afunction of time.

\item Show that solutions to the Neumann problem for th einhomogeneous heat equation
\begin{align*}
	\BINDEF{u_t - \Delta u=f, & x\in U, t>0}{\PART{u}{\nu} = h & x\in \partial U,\\u(x,0)=g(x)}
\end{align*}are unique.  (Hint: sue teh result fo the previous exercise.)

Recall that solutions to Neumann problem for Laplace's equations were not unique.  Why doesn't eh same problem occur here?

\item \begin{enumerate}[label=(\alph*)]

	\item Consider the initial value problem 
	\begin{align*}
		\BINDEF{u_t-\Delta u+b\cdot Du=0, & x\in \R^n, t>0}{u(x,0)=g(x) }
	\end{align*}where $b \in \R^n$ is a fixed vector. (Physically, the $b-Du$ term represents advection.)
	
	Find an explicit soluton to this probem suint eh chang eof variables $z=x+bt$ adn the solution formula for the heat equation (equation (9) in Section 2.3 of Evans).
	
	\item Now consieer a simliar equaiton with asource term:
	\begin{align*}
		xx%BINDEF{u_t-\Delta u+b\cdot Du=f(x,t), & x\in \R^n, t>0}{u(x,0)=g(x)}
	\end{align*}Find an explitic solution to this problem using the same chang eof variables from (a) and formula (17) in Section 2.3 of Evans.
	
\end{enumerate}

\item Consdier the PDE
\begin{align*}
	u_t-\Delta u-cs\cdot Du=0,\, x\in U, t\in (0,T),
\end{align*}where $c>0$ is a positive constant.  Prove that this PDE satisfies a maximum principle, i.e., the maximum of $u$ on $\hat{U}_T$ occurs on the parabolic boundary. (Hint: imitate the proof given in class for th ehet equation on a bounded domain.)

\item the PDE $u_t-xU_{xx}=0$ is a variable-coefficient parabolic equation in one space dimension.  Show that this PDE doe snot satisfy a maximum principle. (Hint:  consider $(u(x,t)=-2xt-x^2$.  Where does the maximum of this function on the closed rectangle $-x\le x\le 2,0\le t\le 1$ occur?)
\end{enumerate}

\end{document}